%CLASSE DOCUMENTO - LINGUA E DIMENSIONE FONT
\documentclass[trieste,12pt]{toptesi}
%%%%%%%%%%%%%%%%%%%%%%%%%%%%%%%%%%%%%%%%%%%%%%%%%%%%%%%%%%%%%%%

% INCLUSIONE PACCHETTI
%\usepackage[utf8x]{inputenc}
\usepackage[utf8]{inputenc} %utf8
\usepackage[italian]{babel}
\usepackage[T1]{fontenc}
\usepackage{blindtext}
\usepackage{graphicx,wrapfig}
\usepackage{booktabs}
\usepackage{lmodern}
\usepackage{varioref}
\usepackage{url}
\usepackage{array}
\usepackage{paralist}{\obeyspaces\global\let =\space}
\usepackage{verbatim} 
\usepackage{subfig}
\usepackage{tabularx}
\usepackage{amsmath}
\usepackage{amsfonts}
\usepackage{float}
\usepackage{amssymb}
\usepackage{multicol}
\usepackage{multirow}
\usepackage{listings}
\usepackage[pass]{geometry}
\usepackage[figuresright]{rotating}
\usepackage{algorithm}
\usepackage{algorithmic}
\usepackage{amsmath}
\usepackage[babel]{csquotes}
\usepackage{hyperref}
\usepackage[backend=bibtex]{biblatex}

%%%%%%%%%%%%%%%%%%%%%%%%%%%%%%%%%%%%%%%%%%%%%%%%%%%%%%%%%%%%%%%

% CONFIGURAZIONE LINK E RIFERIMENTI
\hypersetup{%
    pdfpagemode={UseOutlines},
    bookmarksopen,
    pdfstartview={FitH},
    colorlinks,
    linkcolor={black}, %COLORE DEI RIFERIMENTI AL TESTO
    citecolor={blue}, %COLORE DEI RIFERIMENTI ALLE CITAZIONI
    urlcolor={blue} %COLORI DEGLI URL
}

%%%%%%%%%%%%%%%%%%%%%%%%%%%%%%%%%%%%%%%%%%%%%%%%%%%%%%%%%%%%%%%

% CONFIGURAZIONE LISTATI/CODICE - CANCELLARE SE NON NECESSARIO
% PYTHON - BIANCO E NERO
\lstset{%
	captionpos=b,
	language=Python,
	basicstyle =\small\ttfamily,
	keywordstyle=\color{black}\bfseries,
	breaklines=true,
	breakatwhitespace=true,
	frame=lines,
	numbers=left,
	numberstyle=\footnotesize,
}

%%%%%%%%%%%%%%%%%%%%%%%%%%%%%%%%%%%%%%%%%%%%%%%%%%%%%%%%%%%%%%%

% FRENCHSPACING VA _SEMPRE_ ABILITATO PER DOCUMENTI IN ITALIANO
\frenchspacing

%%%%%%%%%%%%%%%%%%%%%%%%%%%%%%%%%%%%%%%%%%%%%%%%%%%%%%%%%%%%%%%

%DEFINIZIONE SEZIONI IN NUMERAZIONE ROMANA
%ELENCO DEI LISTATI/CODICI
\makeatletter
\newcommand\listofcodes{%
 \iffrontmatter\else\frontmattertrue\fi
 \if@openright\cleardoublepage\else\clearpage\fi
 % change the meaning of \chapter in a group
 \begingroup\def\chapter##1{\@schapter}
 \phantomsection % for the hyperlink
 \lstlistoflistings 
 \endgroup
} 
\makeatother

%%%%%%%%%%%%%%%%%%%%%%%%%%%%%%%%%%%%%%%%%%%%%%%%%%%%%%%%%%%%%%%

% INFORMAZIONI PDF - PERSONALIZZARE
\pdfinfo{%
  /Title    (Riduzione del grado di controversia delle rete sociali connettendo punti di vista opposti)
  /Author   (Stefano Agostini)
  /Subject  (Non definito)
  /Keywords (Controversia)
}

%%%%%%%%%%%%%%%%%%%%%%%%%%%%%%%%%%%%%%%%%%%%%%%%%%%%%%%%%%%%%%%

\ateneo{Università di Roma Tor Vergata}

% LOGO UNIVERSITA
\logosede{images/Logo-Uni-Tor-Vergata.png}

\facolta{Ingegneria}

\corsodilaurea{Ingegneria Informatica}

% TIPOLOGIA TESI
\TesiDiLaurea{Tesi di Laurea Magistrale}

% TITOLO
\titolo{Riduzione del grado di controversia delle reti sociali connettendo punti di vista opposti}

% RELATORE/I
\relatore{Giuseppe F. Italiano}
\secondorelatore{Nikos Parotsidis}

% CANDIDATO - DICITURA (MANTENERE I DUE PUNTI) - CANCELLARE O DECOMMENTARE
%\CandidateName{Candidate:}

% CANDIDATO - NOME E COGNOME
\candidato{\tabular{l}Stefano \textsc{Agostini}\\matricola: 0234240\endtabular}

% DATA - MESE ANNO
\sedutadilaurea{A.A. 2017/2018}

%%%%%%%%%%%%%%%%%%%%%%%%%%%%%%%%%%%%%%%%%%%%%%%%%%%%%%%%%%%%%%%

% LISTA DEI CAPITOLI DA INCLUDERE - PERSONALIZZARE
\includeonly{%
chap_0,%
chap_1,%
chap_2,%
app_a,%
app_b,%
}

% FILE DI BIBLIOGRAFIA
\bibliography{bibliography} 


%%%%%%%%%%%%%%%%%%%%%%%%%%%%%%%%%%%%%%%%%%%%%%%%%%%%%%%%%%%%%%%

% INIZIO DOCUMENTO
\begin{document}

\frontespizio

%%%%%%%%%%%%%%%%%%%%%%%%%%%%%%%%%%%%%%%%%%%%%%%%%%%%%%%%%%%%%%%

%INTERLINEA - DEFAULT 1 - NON ESAGERATE, NON SUPERATE MAI 1.3 ;)
\interlinea{1.3}

%%%%%%%%%%%%%%%%%%%%%%%%%%%%%%%%%%%%%%%%%%%%%%%%%%%%%%%%%%%%%%%

\frontmatter

% DEDICA - PERSONALIZZARE
% VSPACE - PROPORZIONE USATA PER CENTRATURA VERTICALE DEL TESTO
% FLUSHRIGHT - ALLINEAMENTO ORIZZONTALE A DESTRA
\vspace*{\stretch{1}}
\begin{flushright}
\noindent
Da inserire qui tutte le dediche della mia tesi...
\end{flushright}
\vspace*{\stretch{6}}
\cleardoublepage

%%%%%%%%%%%%%%%%%%%%%%%%%%%%%%%%%%%%%%%%%%%%%%%%%%%%%%%%%%%%%%%

% ABSTRACT - PERSONALIZZARE
\sommario
Esistono, e sono sempre esistiti, particolari temi, detti \textit{controversi}, per i quali ognuno di noi si schiera come sostenitore o come oppositore; tali temi possono riguardare contesti politici, sociali o culturali. L'effetto quasi immediato è la divisione della popolazione in due gruppi che hanno visioni opposte sull'argomento \textit{controverso} in considerazione e che difficilmente scambiano tra loro informazioni e punti di vista, non facendo altro che acutizzare la loro inconciliabilità di opinione: gruppi di individui di questo genere sono anche detti \textit{echo chambers}. Gli invidui facenti parte della stessa \textit{echo chamber} hanno quindi le stesse credenze e lo stesso parere riguardo all'argomento \textit{controverso}: essi rinforzano a vicenda le proprie opinioni e sono scarsamente esposti a punti di vista opposti ai propri (i.e. le opinioni che caratterizzano l'altra \textit{echo chamber}). Non è immune da tale fenomeno il mondo dei \textit{social media}. In particolare, \textit{Twitter} identifica un particolare \textit{topic} mediante un \textit{hashtag} (e.g.\textit{\#novax}) e la discussione riguardo a tale \textit{topic} viene descritta mediante un \textit{endorsement graph}, ovvero un \textit{grafo diretto} in cui ciascun nodo rappresenta un utente che partecipa alla discussione e vi è un arco diretto da un nodo \textit{x} ad un nodo \textit{y} se e solo se l'utente \textit{x} approva l'opinione dell'utente \textit{y} (per mezzo dello strumento del \textit{retweet}). Quando il \textit{topic} in considerazione è \textit{controverso}, la struttura dell'\textit{endorsement graph} mette in luce la presenza di gruppi (di utenti) molto connessi al loro interno ma che comunicano poco tra loro: la distanza di opinione che separa tali gruppi (o \textit{echo-chambers}) viene quantificata mediante il così detto \textit{indice di controversia} del grafo, che viene misurato, nel lavoro di tesi proposto, utilizzando una metrica basata sui \textit{random walks}. \textit{Endorsement graphs} con elevati \textit{indici di controversia} sono caratterizzati da \textit{echo chambers} poco connesse tra loro, in ciascuna delle quali è amplificata una visione univoca ed acritica sull'argomento \textit{controverso} considerato.\\Con l'obiettivo di ridurre efficacemente tale \textit{indice di controversia}, è stato implementato un \textit{framework}, utilizzando il linguaggio \textit{Python} e con l'ausilio della libreria \textit{NetworkX}, il quale:

\begin{enumerate}
\item acquisisce i dati necessari per costruire l'\textit{endorsement graph} associato ad un \textit{hashtag controverso}, presente nella \textit{social network} di \textit{Twitter}, fornito in input. Cattura la collezione dei dati necessari mediante l'\textit{API} di \textit{Twitter Tweepy};
\item esegue un algoritmo che identifica le \textit{echo chambers};
\item implementa un \textit{edge-recommendation system} che permette di individuare \textit{k} archi diretti (\textit{k} fornito in \textit{input}) che, se aggiunti al grafo, riducono il suo \textit{grado di controversia};
\item offre un \textit{tool} di visualizzazione degli archi individuati all'interno del grafo.
\end{enumerate}
Con riferimento al punto \textit{3}, il problema di ottimizzazione che si vorrebbe risolvere sarebbe quello di trovare il \textit{set} di \textit{k} archi diretti che, se aggiunti al grafo, minimizzano il \textit{grado di controversia}.
Poiché gli \textit{endorsement graphs} delle \textit{social networks} sono generalmente costituiti da un numero molto elevato di nodi (indicato con \textit{n}), risolvere tale problema di ottimizzazione considerando tutte le possibili combinazioni degli archi a gruppi di \textit{k} (complessità \textit{O}${n^2\choose k}$) (approccio \textit{brute force}) è evidentemente molto costoso dal punto di vista computazionale e molto inefficiente anche per quanto riguarda i tempi di esecuzione. Pertanto l'approccio seguito è quello adottato nell'articolo \cite{garimella:paper}, che consiste nel considerare solo un sottoinsieme degli archi possibili (i.e. un sottoinsieme degli archi non ancora materializzati nell'\textit{endorsement graph}) ed estrarre da questo sottoinsieme i \textit{k} più promettenti. Chiaramente questa soluzione potrebbe restituire archi meno efficaci, per quanto riguarda la riduzione del \textit{grado di controversia} che consentono, rispetto a quelli restituiti dall'approccio \textit{brute force} ma apporta un miglioramento in termini di efficienza computazionale; in particolare, l'euristica che specifica la modalità di scelta del sottoinsieme degli archi candidati è cruciale. L'euristica utilizzata in questo lavoro di tesi è quella proposta nell'articolo \cite{garimella:paper}, la quale
consiste nel considerare solo gli archi diretti che permettono di connettere i vertici di grado alto della prima comunità (o \textit{echo chamber}) con i vertici di grado alto della seconda comunità e viceversa: da questo sottoinsieme di archi vengono estratti i \textit{k} più promettenti in termini di riduzione del \textit{grado di controversia} che consentono.\\La bontà di tale euristica è stata valutata considerando due algoritmi alternativi, utilizzati per estrarre i \textit{k} archi più promettenti dal sottoinsieme considerato: 
\begin{itemize}
\item \textit{non-greedy}: vengono scelti in un solo \textit{step} i \textit{k} archi che porterebbero al \textit{grado di controversia} più basso qualora venissero aggiunti al grafo \textit{individualmente};
\item \textit{greedy}: vengono scelti \textit{k} archi in \textit{k steps}, in ognuno dei quali viene estratto l'arco migliore, tra quelli rimanenti, in termini di decremento del \textit{grado di controversia} che apporterebbe se fosse aggiunto al grafo.
\end{itemize}
Nel lavoro di tesi verranno descritte le modalità di implementazione di tali algoritmi e successivamente verranno confrontati tra loro in termini di efficacia (ossia in termini del decremento del \textit{grado di controversia} che consentono, qualora tutti gli archi che consigliano venissero accettati) ed in termini di tempi di esecuzione ed efficienza computazionale: a tal fine sono stati condotti \textit{tests} su tre \textit{endorsement graphs} di \textit{Twitter} corrispondenti ad \textit{hashtags} particolarmente controversi (\textit{\#beefban,\#russia\_march,\#indiana}), volutamente scelti da contesti sociali e culturali diversi in modo tale da ottenere un'analisi più attendibile.\\L'algoritmo \textit{greedy} si rivelerà più preciso in quanto ad ogni \textit{step} si limita a proporre uno ed un solo arco, ossia l'arco migliore in termini del decremento del \textit{grado di controversia} che consentirebbe se fosse accettato; inoltre, ad ogni \textit{step} dell'algoritmo la scelta dell'arco migliore viene condotta solo dopo aver aggiunto al grafo tutti gli archi consigliati negli \textit{steps} precedenti.\\Al contrario, l'algoritmo \textit{non-greedy} propone in un solo passo i \textit{k} archi migliori utilizzando come metrica il decremento del \textit{grado di controversia} che ciascuno di essi apporterebbe se fosse aggiunto \textit{individualmente}: poichè viene valutato il loro impatto \textit{individuale} e viene ignorato il fatto che tale impatto potrebbe decrementare rispetto a quanto valutato man mano che essi vengono aggiunti al grafo, quest'algoritmo rappresenta un'approssimazione dell'algoritmo \textit{greedy} e consente, in generale, un decremento minore del \textit{grado di controversia}. D'altra parte l'algoritmo \textit{greedy} richiede di scansionare tutti gli archi considerati \textit{k} volte (una volta per \textit{step})  mentre l'algoritmo \textit{non-greedy} una volta sola: questo si traduce in un sostanziale vantaggio dal punto di vista dei tempi di esecuzione, ottenibile utlilizzando l'algoritmo \textit{non-greedy}.\\I test condotti sugli \textit{endorsement graphs} di \textit{Twitter} considerati mostreranno che, spesso, l'algoritmo \textit{non-greedy} conduce a risultati confrontabili con quelli ottenuti dall'algoritmo \textit{greedy} e, pertanto, vista la sua maggiore efficienza computazionale può risultare una scelta più vantaggiosa.
\\Per finire, va sottolineato che, in generale, gli archi che vengono scelti dall'\textit{edge-recommendation system} nella realtà non sempre si materializzano (l'utente potrebbe rigettare il consiglio) e per questo è opportuno considerare come metrica anche la \textit{probabilità di accettazione}. Con il proposito di future estensioni, compresa l'introduzione di tale \textit{probabilità}, il \textit{framework} proposto è implementato in modo da prestarsi perfettamente all'aggiunta di altre metriche per la scelta degli archi\footnote{Nel seguito verrà illustrata l'implementazione dell'algoritmo di \textit{Fagin}, il quale permette di effettuare efficientemente l'ordinamento degli archi considerati secondo il prodotto della loro \textit{probabilità di accettazione} con il decremento dell'\textit{indice di controversia} associato.}.



%%%%%%%%%%%%%%%%%%%%%%%%%%%%%%%%%%%%%%%%%%%%%%%%%%%%%%%%%%%%%%%

% INDICI - ELIMINARE GLI INDICI NON NECESSARI

% INDICE GENERALE
\tableofcontents

% INDICE DELLE FIGURE
\listoffigures

% INDICE DELLE TABELLE
\listoftables

% INDICE DEI CODICI
\listofcodes

%%%%%%%%%%%%%%%%%%%%%%%%%%%%%%%%%%%%%%%%%%%%%%%%%%%%%%%%%%%%%%%

\mainmatter

% INCLUSIONE FILE CAPITOLI - PERSONALIZZARE - TENERE COERENTE CON LISTA IN ALTO
\chapter{Introduzione}
\label{chap:introduzione}

Il fenomeno della polarizzazione degli utenti attorno a \textit{topic controversi} che si propongono nei \textit{social media} è ben noto ed il suo studio è già stato affrontato in alcuni articoli scientifici (tra cui \cite{musco:paper}\cite{morales:paper}). Molto spesso i \textit{social media}, mediante algoritmi di \textit{recommendation}, espongono gli utenti a contenuti che si addicono e sono conformi alle loro opinioni e, pertanto, non fanno altro che aggravare il loro stato di polarizzazione. Il proposito di questa tesi, tuttavia, non è quello di ridurre la polarizzazione dei singoli utenti ma quello di minimizzare il \textit{livello di controversia} dell'intera rete di \textit{endorsement} mediante il \textit{bridging} dei due lati opposti della disputa.\\Il \textit{framework} sviluppato acquisisce da \textit{Twitter} i dati necessari per elaborare \textit{endorsement graphs} di \textit{topic controversi}, ne analizza la struttura estraendone le \textit{echo chambers} ed infine implementa un \textit{edge-recommendation system} che ha lo scopo di ridurre il \textit{grado di controversia} creando connessioni (\textit{bridges}) tra utenti che hanno punti di vista opposti (appartengono a \textit{echo chambers} distinte). Infatti la naturale propensione degli individui a dare credito solo a notizie e contenuti che si addicono al proprio parere fa sì che, in assenza di un intervento esterno di \textit{edge-recommendation}, essi rafforzino sempre più la propria convinzione, anche qualora questa fosse sbagliata o acritica.\\\\Illustriamo ora sinteticamente le nozioni teoriche e le assunzioni grazie alle quali è stato possibile progettare ed implementare il \textit{framework}, partendo dal concetto di \textit{endorsement graph}.
\\Gli \textit{endorsement graphs}, nel particolare ambiente di \textit{Twitter}, prendono anche il nome di \textit{retweet graphs}. Dato un \textit{hashtag} controverso, viene a formarsi naturalmente una discussione a riguardo, nella quale gli utenti esprimono una propria opinione e possono approvare il punto di vista delle così dette \textit{autorità}: nel caso particolare di \textit{Twitter} questa \textit{approvazione} si realizza mediante lo strumento del \textit{retweet}, ossia se l'utente \textit{u} fa \textit{retweet} di un \textit{tweet} prodotto dall'utente \textit{v} allora ne approva l'opinione. Ne deriva la formazione di un grafo diretto \textit{G(V,E)} costituito da \textit{n} nodi (gli utenti che partecipano alla discussione) ed i cui archi (\textit{retweets}) esprimono relazioni di condivisione di opinione: data questa struttura dei \textit{retweet graphs}, in caso di discussioni su \textit{topics} (\textit{hashtags}) particolarmente controversi non è sorprendente che esistano \textit{echo chambers}.\\Le \textit{echo-chambers} sono due sottoinsiemi dei nodi \textit{X,Y}, ben separati tra loro (vi sono pochi archi che li congiungono) e tali che \textit{$X \cup Y = V$} e \textit{$X \cap Y =  \emptyset$}. Tale ripartizione dei nodi può essere ottenuta mediante l'utilizzo di un algoritmo di \textit{graph-partitioning}. Nel lavoro proposto è stato utilizzato l'algoritmo di \textit{Girvan-Newman}: esso è un metodo gerarchico usato per rilevare le comunità in sistemi complessi e la cui esecuzione produce un \textit{dendrogramma}, le cui foglie sono i nodi del grafo. 
\\L'\textit{edge-recommendation system} sviluppato utilizza una metrica basata sul concetto di \textit{random walk} per misurare il \textit{grado di controversia} associato al \textit{topic} analizzato (attorno al quale si svolge la discussione nella \textit{social network} di \textit{Twitter}). Più precisamente, per misurare il \textit{grado di controversia} della rete, viene utilizzata la funzione \textit{Random-walk controversy score}:
\\\\
$\textit{RWC(G,X,Y) = (c\textsubscript{x} - c\textsubscript{y})\textsuperscript{T}(r\textsubscript{x} - r\textsubscript{y})}$
\\\\Dove \textit{c\textsubscript{x}} è un vettore di dimensione \textit{n} (numero di vertici dell'\textit{endorsement graph}) che ha valore \textit{1} nelle coordinate corrispondenti ai vertici di grado alto dell'\textit{echo-chamber X} e \textit{0} altrove; similmente viene definito \textit{c\textsubscript{y}}. Infine \textit{r\textsubscript{x}} è il vettore di \textit{PageRank} personalizzato per un \textit{random walk} che parte dai nodi dell'\textit{echo-chamber X}; similmente viene definito \textit{r\textsubscript{y}}. Valori alti di \textit{RWC(G,X,Y)} indicano che, all'equilibrio del \textit{random walk}, è bassa la probabilità di essere nell'\textit{echo-chamber} opposta a quella di partenza: questo è indice di elevata \textit{controversia}. 
\\Il sistema implementato permette, utilizzando in alternativa un algoritmo \textit{greedy} o uno \textit{non-greedy}, di proporre \textit{k} archi (che farebbero da \textit{bridges} tra le due \textit{echo chambers}) in grado di ridurre questo indice; in particolare, l'\textit{edge-recommendation system} proposto restituisce archi la cui efficacia approssima quella degli archi che sono soluzione del problema di ottimizzazione originario che, come detto precedentemente, ha una complessità di un livello troppo elevato (\textit{O}${n^2\choose k}$) per essere risolto in tempi accettabili. Di seguito la definizione del problema di ottimizzazione originario:
\\
\begin{equation*}
\begin{aligned}
& \underset{E\textsubscript{k}}{\text{minimize}}
& & RWC(G(V, E \cup E\textsubscript{k}),X,Y) \\
& \text{subject to}
& & E \cap E\textsubscript{k} = \emptyset, \left|{E\textsubscript{k}}\right| = k
\end{aligned}
\end{equation*}
\\
L'euristica proposta in questo lavoro permette di restringere il dominio degli archi considerati, consentendo di ottenere risultati paragonabili con quelli ottenibili mediante la soluzione del problema di ottimizzazione appena descritto e con il vantaggio di avere una complessità computazionale di molto inferiore; essa, forniti in \textit{input} i valori di \textit{k\textsubscript{1}} e \textit{k\textsubscript{2}}:
\begin{enumerate}
\item considera i \textit{k\textsubscript{1}} vertici con \textit{in-degree} più alto dell'\textit{echo-chamber X} e i \textit{k\textsubscript{2}} vertici con \textit{in-degree} più alto dell'\textit{echo-chamber Y};
\item costruisce il dominio degli archi considerati come l'insieme di tutti i possibili archi diretti, non presenti ancora nel grafo, che abbiano come estremi un vertice dell'insieme \textit{K\textsubscript{1}} e uno dell'insieme \textit{K\textsubscript{2}};  
\end{enumerate}
Gli algoritmi \textit{greedy} e \textit{non-greedy} operano entrambi a partire dal dominio così definito. Hanno lo scopo di proporre i \textit{k} archi migliori del dominio in termini del decremento del grado di controversia \textit{RWC(G,X,Y)} che consentirebbero qualora si materializzassero nel grafo; la differenza sta nella modalità di selezione. In particolare:
\begin{enumerate}
\item l'algoritmo \textit{greedy} sceglie \textit{avidamente} ognuno dei \textit{k} archi: esso impiega \textit{k} passi, in ognuno dei quali sceglie l'arco più appetibile del dominio (ossia l'arco del dominio associato al $\delta RWC$ maggiore) e lo aggiunge al grafo.
\item
\end{enumerate}
%qui descrivo sinteticamente i due algoritmi
In pratica, per ridurre la \textit{controversia}, si propone ad un certo utente il contenuto (i.e. un \textit{tweet}) di un altro utente che ha posizioni opposte rispetto alle sue sull'argomento, sperando che possa accettarne il punto di vista mediante un \textit{retweet}: ciò provocherebbe la formazione di un arco tra le due comunità (\textit{echo chambers}) con l'effetto di ridurre il \textit{grado di controversia} dell'intero \textit{retweet graph}.  

%\section{Quaraquaraqua}
%\label{sec:quaqaraqua}

%This is a reference to a chapter \ref{chap:quo}. This is a reference to a figure \ref{fig:doge}. This is a reference to some code \ref{lst:hello}. This is a citation %\cite{famous:paper}.

%\lstinputlisting[basicstyle=\fontsize{8}{10}\selectfont\ttfamily ,label=lst:prova, firstline=1, lastline=34, caption={I directly included a portion of a file}]{code/prova.py}

%\begin{lstlisting}[language=Java, label=lst:java, caption={Some code in another language than the default one}]
%public void prepare(AClass foo) {
%        AnotherClass bar = new AnotherClass(foo)
%}
%\end{lstlisting}

%\Blindtext

%\begin{figure}
%\begin{center}
%\includegraphics[width=0.5\columnwidth]{images/doge.png}
%\end{center}
%\caption{This is not a figure. It's a caption.}
%\label{fig:doge}
%\end{figure}

\chapter{Lo stato dell'arte}
\label{chap:stato}
Nonostante il \textit{Web} offra metodi immediati per accedere a qualsiasi tipo di informazione, molti studi (e.g. \textit{Liao et al.}\cite{liao:paper}) dimostrano che, quando un utente può scegliere, esso preferisce essere esposto a contenuti che non contrastano con la sua opinione. Questo fenomeno ha portato ad una maggiore frammentazione e polarizzazione online.

\chapter{Teoria alla base del problema ed algoritmi per la risoluzione}
\label{chap:teoria}
\section{Misura del grado di controversia}
Prima di dare una definizione formale del \textit{random-walk controversy score}, elenchiamo ed illustriamo i passi necessari per calcolarlo.
\begin{enumerate}
\item Fissato il \textit{topic t} per il quale si vuole quantificare il grado di controversia, è possibile descrivere la discussione mediante l'\textit{endorsement graph G(V,E)}. Nell'ambiente di \textit{Twitter}, il \textit{topic t} è identificato da un \textit{hashtag} (e.g.\textit{\#hashtag}) ed i nodi del grafo rappresentano gli utenti che hanno preso parte alla discussione utilizzando almeno una volta tale \textit{hashtag} nei loro \textit{tweets}; gli archi del grafo identificano i \textit{retweets} tra gli utenti, che esprimono relazioni di condivisione di opinione riguardo al \textit{topic}.
\item Ipotizzando che il \textit{topic t} sia controverso, è possibile partizionare i nodi del grafo \textit{G(V,E)} in due insiemi \textit{X},\textit{Y} ben separati tra loro (i.e. vi sono pochi archi che li interconnettono). Tali insiemi quindi soddisfano le seguenti proprietà:
\begin{enumerate}
\item $\textit{X}\cup\textit{Y} = \textit{V}$;
\item $\textit{X}\cap\textit{Y} = \emptyset$.
\end{enumerate}
Gli insiemi \textit{X} ed \textit{Y} rappresentano i due lati della controversia (i.e. le \textit{echo-chambers}).
\\Per identificare le \textit{echo-chambers}, nell'implementazione proposta è stato utilizzato l'algoritmo di \textit{graph-partitioning} di \textit{Girvan-Newman}\cite{girvan:paper}. Tale algoritmo agisce rimuovendo progressivamente archi dal grafo originario: l'esecuzione viene arrestata quando la rimozione degli archi ha portato ad individuare due comunità distinte che non comunicano (i.e. non sono collegate da nessun arco). La metrica utilizzata da \textit{Girvan-Newman} per identificare l'arco da rimuovere ad ogni passo è la così detta \textit{edge-betweenness centrality}: dato un arco \textit{e}, essa è definita come \textit{il numero di cammini di costo minimo tra coppie di nodi del grafo che passano attraverso l'arco e}. Nel caso in cui vi sia più di un percorso di costo minimo tra una coppia di nodi, a ciascun percorso viene assegnato uguale peso in modo tale che il peso totale di tutti i percorsi sia uguale all'unità. Di seguito la formula che definisce questa metrica di centralità:
\\
\begin{equation}b(e) = \sum_{\substack{
   s \neq t
  }} 
 \cfrac{\sigma\textsubscript{st}(e)}{\sigma\textsubscript{st}}
\end{equation}
\\ 
Dove \textit{$\sigma\textsubscript{st}$} è il numero totale di percorsi di costo minimo dal nodo \textit{s} al nodo \textit{t} e \textit{$\sigma\textsubscript{st}(e)$} è il numero di tali percorsi che passano attraverso l'arco \textit{e}. 
\\L'intuizione è: se la struttura del grafo è caratterizzata da due comunità di nodi connesse tra loro da pochissimi archi, allora tutti i percorsi tra queste due comunità dovranno passare attraverso questi archi. Ne consegue che quest'ultimi saranno caratterizzati da un'alta \textit{betweenness centrality}. Sfruttando la peculiarità di tali archi, l'algoritmo di \textit{Girvan-Newman} si rivela un ottimo metodo per rilevare le \textit{echo-chambers}.
\item A questo punto è possibile procedere con la definizione del \textit{random-walk controversy score}. L'\textit{RWC} è definito come \textit{la differenza della probabilità che un random walk che parte da una echo-chamber all'equilibrio vi permanga e la probabilità che invece tale random walk all'equilibrio finisca nell'echo-chamber opposta}. Tale misura viene calcolata mediante l'utilizzo di due esecuzioni dell'algoritmo di \textit{PageRank} personalizzato, le quali non sono altro che due \textit{random walks} particolari.\\
\textit{PageRank} è un algoritmo di analisi che assegna un peso numerico a ciascun nodo di un grafo diretto, con lo scopo di quantificare la sua importanza relativa. Le applicazioni più frequenti di \textit{PageRank} riguardano l'ambito del \textit{World Wide Web}, in cui i grafi hanno come nodi le pagine \textit{web} e come archi i collegamenti ipertestuali. Ciò non toglie che \textit{PageRank} sia uno strumento molto potente anche nell'ambito delle reti sociali, poiché riesce a quantificare l'importanza di utente nell'ambito di una discussione: tale importanza misura il suo grado di popolarità e di rilevanza.
\\Per il calcolo dell'\textit{RWC} vengono utilizzate due esecuzioni distinte dell'algoritmo di \textit{PageRank}, indicate con \textit{page\textsubscript{x} e page\textsubscript{y}}, ognuna delle quali opera sul \textit{retweet graph} corrispondente al \textit{topic t} in input ma inizia il suo \textit{random-walk} partendo, rispettivamente, da uno dei nodi della \textit{comunità X} e da uno dei nodi della \textit{comunità Y} (\textit{comunità = echo-chamber}). Inoltre, \textit{page\textsubscript{x} e page\textsubscript{y}} ad ogni passo possono decidere di continuare il \textit{random-walk} (potendo scegliere con uguale probabilità uno degli archi in uscita dal nodo in cui si trovano attualmente) o di ricominciare il proprio cammino (\textit{restart}), tornando, rispettivamente, in uno dei nodi della \textit{comunità X} ed in uno dei nodi della \textit{comunità Y}: la seconda scelta viene compiuta con una probabilità detta di \textit{restart} e, chiaramente, la prima con una probabilità che ne è il complementare (i.e. la somma delle probabilità deve restituire \textit{1}).
\\Scendendo più nel dettaglio, siano:
\begin{itemize}
\item \textit{P} la \textit{matrice delle probabilità di transizione per colonna}\footnote{Se il \textit{retweet graph} in considerazione ha \textit{N} nodi, tale matrice ha dimensione $N \times N$ ed ogni suo elemento \textit{P[i][j]} è la probabilità di passare dal nodo \textit{j} al nodo \textit{i} in un solo passo, sapendo di essere attualmente nel nodo \textit{j}.} associata al \textit{retweet graph} considerato;
\item \textit{X\textsuperscript{*} e Y\textsuperscript{*}} rispettivamente gli insiemi dei \textit{k\textsubscript{1} e k\textsubscript{2}} nodi con \textit{in-degree} più alto delle due comunità \textit{X e Y}. Inoltre sia \textit{c\textsubscript{x}} un vettore di dimensione \textit{n} avente valore \textit{1} nelle coordinate corrispondenti ai nodi dell'insieme \textit{X\textsuperscript{*}} e \textit{0} altrove; similmente viene definito \textit{c\textsubscript{y}};
\item \textit{r\textsubscript{x}} il vettore di \textit{PageRank} personalizzato per il \textit{random walk} che parte dalla comunità \textit{X}. Sia inoltre $(1-\alpha)$ la probabilità di \textit{restart} di tale \textit{random walk} (e dunque $\alpha$ è la probabilità di continuare) e sia $\textit{e\textsubscript{x} = Uniform(X)}$ il suo vettore di \textit{restart}: ossia il \textit{random walk}, ad ogni passo, decide di ricominciare il proprio cammino con probabilità $(1-\alpha)$ e tra tutti i nodi della \textit{comunità X}, con eguale probabilità, sceglie il nodo da cui ricominciare.
\\Simili considerazioni valgono per \textit{r\textsubscript{y}}.
\end{itemize}
Bisogna ora risolvere il problema dei vertici \textit{dangling}, ossia i vertici del grafo che non hanno archi in uscita. Se un \textit{random walk} dovesse casualmente finire in uno di questi nodi esso potrebbe non uscirne, compromettendo l'esecuzione di \textit{PageRank}. Per evitare tale situazione, vengono utilizzate non una ma due \textit{matrici delle probabilità di transizione per colonna P\textsubscript{x} e P\textsubscript{y}}, usate rispettivamente dal \textit{random walk} che inizia dalla comunità \textit{X} e dal \textit{random walk} che inizia dalla comunità \textit{Y}. Se il grafo non contiene vertici \textit{dangling}, si ha banalmente \textit{P\textsubscript{x} = P\textsubscript{y} = P}; se al contrario li contiene, le matrici \textit{P\textsubscript{x} e P\textsubscript{y}} sono definite in modo tale che le probabilità di transizione dai vertici \textit{dangling} sono uguali, rispettivamente, ai vettori di \textit{restart} \textit{e\textsubscript{x} ed e\textsubscript{y}}.
\\Il \textit{PageRank} personalizzato per i due \textit{random walks} che, rispettivamente, iniziano nella \textit{comunità X e Y} e dato da:
\\
\begin{equation}
r_x = \alpha P_xr_x + (1-\alpha)e_x
\end{equation}
\begin{equation}
r_y = \alpha P_yr_y + (1-\alpha)e_y
\end{equation}
\\
Per il calcolo dei vettori di \textit{PageRank} personalizzati (\textit{r\textsubscript{x} e r\textsubscript{y}}), quindi, bisogna imporre una situazione di \textit{stazionarietà} dei rispettivi \textit{random walks}. 
\\Possiamo finalmente definire il \textit{random-walk controversy score}:
\\
\begin{equation}
RWC(G,X,Y) = (c_x - c_y)\textsuperscript{T}(r_x - r_y)
\end{equation}
\\
Sostituendo le equazioni (3.2) e (3.3) nell'espressione (3.4) si ottiene:
\\
\begin{equation}
RWC(G,X,Y) = (1 - \alpha)(c_x - c_y)\textsuperscript{T}(M_x\textsuperscript{-1}e_x - M_y\textsuperscript{-1}e_y)
\end{equation}
\\
Dove \textit{M\textsubscript{x} = (I - $\alpha$P\textsubscript{x})} e \textit{M\textsubscript{y} = (I - $\alpha$P\textsubscript{y})}. Quest'ultima è la formula utilizzata nell'implementazione del sistema proposto. Ad ogni modo, riferendosi all'equazione (3.4) è possibile fare le seguenti osservazioni:
\begin{enumerate}
\item Innanzitutto occorre precisare che \textit{r\textsubscript{x}} è una \textit{distribuzione stazionaria di probabilità} del \textit{random walk} associato all'esecuzione di \textit{PageRank page\textsubscript{x}}: ognuno dei suoi elementi \textit{r\textsubscript{x}[i]} ($\forall \textit{i = 0,...,n-1}$) rappresenta la probabilità che il \textit{random walk} si trovi nel nodo \textit{i} all'equilibrio. Simili considerazioni valgono per \textit{r\textsubscript{y}};
\item Qualora le due comunità (\textit{echo-chambers}) \textit{X e Y} fossero molto divise tra loro (i.e. pochissimi archi diretti che le connettono), ci si aspetterebbe un \textit{r\textsubscript{x}} con valori pressoché nulli nelle coordinate corrispondenti ai nodi di \textit{Y} e non nulli altrove e, viceversa, parlando di \textit{r\textsubscript{y}}\footnote{Ricordare che, essendo \textit{distribuzioni di probabilità}, vale: $\sum_{i=0}^{n-1} r_x[i] = 1$ e $\sum_{j=0}^{n-1} r_y[j] = 1$.}. In tal caso l'espressione $(c_x - c_y)\textsuperscript{T}(r_x - r_y)$ assume il valore massimo, indice di un'elevata \textit{controversia} nella rete.
\item Qualora le due comunità fossero invece sufficientemente connesse tra loro, ci si aspetterebbe un vettore \textit{r\textsubscript{x}} con valori abbastanza uniformi su tutte le sue coordinate; la stessa considerazione varrebbe per \textit{r\textsubscript{y}}. In tal caso il grafo non presenterebbe un'elevata controversia (\textit{RWC(G,X,Y)} assume valori bassi), e questo risultato sarebbe confermato dal fatto che le due comunità sono sufficientemente esposte l'una a l'altra.
\end{enumerate}
\end{enumerate}
Per terminare, si può affermare che l'indice \textit{RWC(G,X,Y)} riesce a descrivere perfettamente il grado di controversia della discussione in atto, in funzione del livello di esposizione reciproca delle due comunità (ossia i due punti di vista opposti): qualora tale livello fosse molto basso, l'indice di controversia sarebbe molto elevato e, pertanto, la scelta più efficace per attenuarlo sarebbe quella di cercare di esporre gli utenti a visioni opposte alle proprie. 
\\Nel prossimo paragrafo saranno illustrati gli algoritmi utilizzati nel sistema proposto, i quali hanno come obiettivo proprio quello di esporre reciprocamente, \textit{nel modo più efficace possibile}, i due lati della controversia. 

\section{Definizione formale degli algoritmi per la risoluzione}
Il problema che si vorrebbe risolvere è il seguente: 
\\\\
\textit{Dato un endorsement graph G(V,E) che descrive una discussione riguardo ad un certo topic controverso t, trovare l'insieme di k archi diretti non ancora presenti nel grafo che, qualora si materializzassero, minimizzerebbero il suo RWC(G,X,Y).}
\\\\
Ovvero, formalmente:
\\
\begin{equation*}
\begin{aligned}
& \underset{E\textsubscript{k}}{\text{minimize}}
& & RWC(G(V, E \cup E\textsubscript{k}),X,Y) \\
& \text{subject to}
& & E\textsubscript{k} \subseteq V \times V \setminus E, \left|{E\textsubscript{k}}\right| = k
\end{aligned}
\end{equation*}
\\
Risolvere tale problema, così come si presenta, richiederebbe di considerare tutti le possibili combinazioni degli archi ancora non presenti nel grafo presi a gruppi di \textit{k}: queste combinazioni sono pari a \textit{O}($n^2 \choose k$), dove con \textit{n} indichiamo il numero di nodi del grafo.
\\Come è facile immaginare, un algoritmo che considera tutte queste combinazioni di archi è molto inefficiente dal punto di vista computazionale, soprattutto alla luce del fatto che gli \textit{endorsement graphs} delle reti sociali come quella di \textit{Twitter} sono costituiti da alcune migliaia di nodi e migliaia di archi.
\\Si propone, pertanto, di restringere il dominio degli archi candidati: seguendo l'approccio dell'articolo \cite{garimella:paper} è possibile considerare solo gli archi tra vertici con \textit{in-degree} alto di ciascuna \textit{echo-chamber}. La scelta di questa euristica, come osservano gli autori dell'articolo \cite{garimella:paper}, è giustificata dal fatto che la struttura degli \textit{endorsement graph} spesso mette in luce la presenza di un piccolo numero di nodi \textit{leader} (i.e. utenti popolari) ed un gran numero di nodi \textit{non leader} (i.e. utenti seguaci): i nodi \textit{leader} ricevono molti archi in ingresso (ossia hanno alto \textit{in-degree}) in quanto i loro \textit{tweet} vengono \textit{retweettati} da molti nodi \textit{non leader} (ossia ricevono molta \textit{approvazione} dai loro seguaci, i quali si fidano della loro opinione). Intuitivamente, qualora si riuscisse a convincere i nodi \textit{leader} ad approvare contenuti che esprimono opinioni sul \textit{topic controverso} opposte alla propria, i nodi \textit{non leader} (i seguaci) sarebbero spinti a fare altrettanto.
Questa osservazione suggerisce che gli archi tra vertici con \textit{in-degree} alto di ciascuna \textit{echo-chamber} sono buoni candidati per ottenere un abbassamento del \textit{random-walk controversy score}\footnote{Si rimanda alla lettura del \textit{Teorema 1.} dell'articolo \cite{garimella:paper}.}. 
\\Mediante questo ridimensionamento del dominio degli archi candidati, ci si propone di ridurre l'\textit{RWC} esponendo alcuni nodi \textit{leader} di ciascuna \textit{echo-chamber} ai contenuti di alcuni nodi \textit{leader} dell'\textit{echo-chamber} opposta. 
\\Ora, con l'obiettivo di effettuare la scelta dei \textit{k} archi migliori appartenenti al dominio così definito, vengono proposti due algoritmi: un \textit{algoritmo non greedy} (indicato di seguito come \textit{Algorithm 1}) ed un \textit{algoritmo greedy}.
\\L'algoritmo \textit{non greedy} ha un tempo di esecuzione pari a \textit{O($k_1 \cdot k_2$)}, che costituisce un ottimo \textit{speedup} rispetto all'algoritmo \textit{brute-force}, che considera invece tutte le combinazioni di archi. \textit{Algorithm 1} si limita a scegliere dal dominio degli archi, ridimensionato come appena detto, i \textit{k} archi migliori in termini del decremento dell'\textit{RWC} che ciascuno consente se aggiunto \textit{individualmente}. Tuttavia il \textit{decremento dell'RWC} che un qualsiasi arco \textit{"e"} consente individualmente (\textit{$\delta$RWC\textsubscript{e}}) si rivelerebbe tale solo se tale arco fosse aggiunto per primo al grafo; più precisamente $\forall E' \supseteq E, \forall e \in V \times V \setminus E'$ vale l'espressione (3.6). Ovvero la somma dei \textit{$\delta$RWC} dei \textit{k} archi scelti dall'algoritmo \textit{non greedy} non corrisponde al decremento \textit{reale} dell'\textit{RWC} che si osserverebbe se tali archi apparissero nel grafo, ma ne è un \textit{upper-bound}.

\begin{equation}
\resizebox{1 \textwidth}{!} 
{
    $ RWC(G(V,E),X,Y) - RWC(G(V, E \cup \{e\}),X,Y) \geq RWC(G(V,E'),X,Y) - RWC(G(V, E' \cup \{e\}),X,Y)$
} 
\end{equation}

\begin{algorithm}
\caption{Algoritmo \textit{non greedy} per la scelta dei \textit{k} archi}
\begin{algorithmic} 
\REQUIRE Il grafo \textit{G} e le comunità \textit{X,Y}; il numero di archi da proporre \textit{k}; i \textit{k\textsubscript{1} e k\textsubscript{2}} vertici con \textit{in-degree} più alto in \textit{X e Y}, rispettivamente 
\ENSURE La lista dei \textit{k} archi migliori, in termini del decremento dell'\textit{RWC} che consentono se aggiunti individualmente
\STATE Initialization : Output $\leftarrow$ lista vuota;
\FOR{\textit{i = 1:k\textsubscript{1}}} 
\STATE nodo $u = X[i]$;
\FOR{\textit{j = 1:k\textsubscript{2}}}
\STATE nodo $v = Y[j]$;
\STATE Calcola il decremento dell'\textit{RWC $\delta$RWC\textsubscript{u $\rightarrow$ v}} che si otterrebbe qualora l'arco \textit{(u,v)} venisse aggiunto al grafo;
\STATE Aggiungi \textit{$\delta$RWC\textsubscript{u $\rightarrow$ v}} alla lista \textit{Output};
\STATE Calcola il decremento dell'\textit{RWC $\delta$RWC\textsubscript{v $\rightarrow$ u}} che si otterrebbe qualora l'arco \textit{(v,u)} venisse aggiunto al grafo;
\STATE Aggiungi \textit{$\delta$RWC\textsubscript{v $\rightarrow$ u}} alla lista \textit{Output};
\ENDFOR
\ENDFOR
\STATE Lista \textit{Output} ordinata $\leftarrow$ Ordina la lista \textit{Output} secondo i \textit{$\delta$RWC} (presi in modulo) in ordine decrescente;
\RETURN I migliori \textit{k} dalla lista \textit{Output} ordinata;
\end{algorithmic}
\end{algorithm}

\section{Calcolo del decremento della controversia associato ad un arco}


\appendix
% INCLUSIONE APPENDICI - - PERSONALIZZARE - TENERE COERENTE CON LISTA IN ALTO
\chapter{An appendix}
\label{app:a}
\Blindtext

\chapter{Another appendix}
\label{app:b}
\Blindtext

%%%%%%%%%%%%%%%%%%%%%%%%%%%%%%%%%%%%%%%%%%%%%%%%%%%%%%%%%%%%%%%

% BIBLIOGRAFIA
\addcontentsline{toc}{chapter}{\refname}
\nocite{*}
\printbibliography

\end{document}
